% GENERATED FILE. Edit canonical lessons, then run npm run generate:books.
% canonical-source-hash: fnv1a64:5eff68393958bbb9
% canonical-lessons: GU-C136-khali, GU-C136-bharelu, GU-C136-sacu, GU-C136-khotu, GU-C136-amir

\chapter{Empty, Full, True, False, Rich}
\label{ch:gu-empty-full-true-false-rich}
\hlchaptermodality{136}

\begin{chapteropening}
\textbf{By the end of this chapter:} \emph{I can say empty, full, true, false and rich.}

Complete the last lesson of chapter 136: i can say empty, full, true, false and rich.
\end{chapteropening}

\section[\texorpdfstring{\gu{ખાલી} (khālī)}{ખાલી (khālī)}]{\gu{ખાલી} (khālī) — empty}

\label{lesson:GU-C136-khali}



\begin{quote}
Before the new one: say the Gujarati for ripe, then the Gujarati for raw, unripe.
\end{quote}

\subsection*{You'll want to know: \gu{ખાલી}}
\textbf{\gu{ખાલી}} (\emph{khālī}) — \textquotedblleft{}empty\textquotedblright{}.

\begin{grammarlens}[title={What you've built}]
One more word for this chapter.
\end{grammarlens}

\subsection*{Your turn}
\begin{itemize}
\raggedright
  \item \emph{Say it:} \emph{khālī}
  \item \emph{Say it:} \emph{khālī}, once more
  \item \emph{Say it:} say \emph{khālī}
  \item \emph{Recall:} say the Gujarati for ripe, then the Gujarati for raw, unripe, then say \emph{khālī} again
\end{itemize}

\begin{tcolorbox}[breakable,colback=teal!4,colframe=teal!35!black,title={Before you move on}]
What does \gu{ખાલી} mean? (\textquotedblleft{}Empty\textquotedblright{}.) Say it once more.
\end{tcolorbox}
\section[\texorpdfstring{\gu{ભરેલું} (bharelũ)}{ભરેલું (bharelũ)}]{\gu{ભરેલું} (bharelũ) — full}

\label{lesson:GU-C136-bharelu}



\begin{quote}
Before the new one: say the Gujarati for raw, unripe, then the Gujarati for empty.
\end{quote}

\subsection*{You'll want to know: \gu{ભરેલું}}
\textbf{\gu{ભરેલું}} (\emph{bharelũ}) — \textquotedblleft{}full\textquotedblright{}.

\begin{grammarlens}[title={What you've built}]
One more word for this chapter.
\end{grammarlens}

\subsection*{Your turn}
\begin{itemize}
\raggedright
  \item \emph{Say it:} \emph{bharelũ}
  \item \emph{Say it:} \emph{bharelũ}, once more
  \item \emph{Say it:} say \emph{bharelũ}
  \item \emph{Recall:} say the Gujarati for raw, unripe, then the Gujarati for empty, then say \emph{bharelũ} again
\end{itemize}

\begin{tcolorbox}[breakable,colback=teal!4,colframe=teal!35!black,title={Before you move on}]
What does \gu{ભરેલું} mean? (\textquotedblleft{}Full\textquotedblright{}.) Say it once more.
\end{tcolorbox}
\section[\texorpdfstring{\gu{સાચું} (sācũ)}{સાચું (sācũ)}]{\gu{સાચું} (sācũ) — true, right}

\label{lesson:GU-C136-sacu}



\begin{quote}
Before the new one: say the Gujarati for empty, then the Gujarati for full.
\end{quote}

\subsection*{You'll want to know: \gu{સાચું}}
\textbf{\gu{સાચું}} (\emph{sācũ}) — \textquotedblleft{}true, right\textquotedblright{}.

\begin{grammarlens}[title={What you've built}]
One more word for this chapter.
\end{grammarlens}

\subsection*{Your turn}
\begin{itemize}
\raggedright
  \item \emph{Say it:} \emph{sācũ}
  \item \emph{Say it:} \emph{sācũ}, once more
  \item \emph{Say it:} say \emph{sācũ}
  \item \emph{Recall:} say the Gujarati for empty, then the Gujarati for full, then say \emph{sācũ} again
\end{itemize}

\begin{tcolorbox}[breakable,colback=teal!4,colframe=teal!35!black,title={Before you move on}]
What does \gu{સાચું} mean? (\textquotedblleft{}True, right\textquotedblright{}.) Say it once more.
\end{tcolorbox}
\section[\texorpdfstring{\gu{ખોટું} (khoṭũ)}{ખોટું (khoṭũ)}]{\gu{ખોટું} (khoṭũ) — false, wrong}

\label{lesson:GU-C136-khotu}



\begin{quote}
Before the new one: say the Gujarati for full, then the Gujarati for true.
\end{quote}

\subsection*{You'll want to know: \gu{ખોટું}}
\textbf{\gu{ખોટું}} (\emph{khoṭũ}) — \textquotedblleft{}false, wrong\textquotedblright{}.

\begin{grammarlens}[title={What you've built}]
One more word for this chapter.
\end{grammarlens}

\subsection*{Your turn}
\begin{itemize}
\raggedright
  \item \emph{Say it:} \emph{khoṭũ}
  \item \emph{Say it:} \emph{khoṭũ}, once more
  \item \emph{Say it:} say \emph{khoṭũ}
  \item \emph{Recall:} say the Gujarati for full, then the Gujarati for true, then say \emph{khoṭũ} again
\end{itemize}

\begin{tcolorbox}[breakable,colback=teal!4,colframe=teal!35!black,title={Before you move on}]
What does \gu{ખોટું} mean? (\textquotedblleft{}False, wrong\textquotedblright{}.) Say it once more.
\end{tcolorbox}
\section[\texorpdfstring{\gu{અમીર} (amīr)}{અમીર (amīr)}]{\gu{અમીર} (amīr) — rich}

\label{lesson:GU-C136-amir}



\begin{quote}
Before the new one: say the Gujarati for true, then the Gujarati for false.
\end{quote}

\subsection*{You'll want to know: \gu{અમીર}}
\textbf{\gu{અમીર}} (\emph{amīr}) — \textquotedblleft{}rich\textquotedblright{}.

\begin{grammarlens}[title={What you've built}]
One more word for this chapter.
\end{grammarlens}

\subsection*{Your turn}
\begin{itemize}
\raggedright
  \item \emph{Say it:} \emph{amīr}
  \item \emph{Say it:} \emph{amīr}, once more
  \item \emph{Say it:} say \emph{amīr}
  \item \emph{Recall:} say the Gujarati for true, then the Gujarati for false, then say \emph{amīr} again
\end{itemize}

\begin{tcolorbox}[breakable,colback=teal!4,colframe=teal!35!black,title={Before you move on}]
What does \gu{અમીર} mean? (\textquotedblleft{}Rich\textquotedblright{}.) Say it once more.
\end{tcolorbox}
